% \documentclass[12pt]{report}
% \usepackage[table]{xcolor}
% \usepackage[margin=1in]{geometry}
% \usepackage{booktabs}
% \usepackage{hyperref}
% \usepackage{array}
% \usepackage{float}
% \usepackage{caption}
% \usepackage{longtable}
% \usepackage{booktabs}     % for \toprule, \midrule, \bottomrule
% \usepackage{caption}      % nicer caption handling

% \geometry{margin=1in}
% \begin{document}

% \chapter*{Document 15: Conclusions on Designed Product}

% % 15.1 Summary of Design Project Results
% \section*{15.1 Summary of Design Project Results}
% This design project, undertaken in collaboration with the USC Office of Sustainability and USC Facilities Department, focused on retrofitting and enhancing the existing 1,000‑gallon rainwater collection tank at the Green Quad permaculture garden—a 0.5‑acre, student-run organic garden behind a LEED‑certified building. The tank’s original flow of 1–2\,gpm was insufficient for operational needs.

% \textbf{Objectives:} Increase flow rate for rapid filling of watering cans; enable sprinkler operation; ensure a simple, quiet, CO$_2$‑free system; integrate aesthetically into the garden; maximize lifespan with zero maintenance; and maintain reliability under failure modes.
% \newline\newline
% The final solution implements a two‑tier architecture:
% \begin{itemize}
%   \item \textbf{Passive Gravity‑Driven System:} Brass valve with extended lever delivers over 20 gpm under peak head (900+ gallons in tank) and operates with as little as 200 gallons, filling a standard watering can in 30 seconds at reduced head conditions.
%   \item \textbf{Active Solar‑Powered Pressurized System:} Solar‑charged battery powers a UL‑listed pump maintaining 12–16 gpm at 20–40 psi, supports sprinklers via dual spigots, and automatically switches to mains power backup if necessary.
% \end{itemize}





% % 15.3 Recommendations for Future Development
% % 15.3 Recommendations for Future Development
% \section*{15.3 Recommendations for Future Development}
% \begin{itemize}
%   \item \textbf{Readiness:} Prototype is functional; commercial deployment requires durability testing and cost refinement.
%   \item \textbf{Design Refinements:} Enhance material selection for longevity; standardize modular connectors; refine brass lever geometry for ergonomic operation.
%   \item \textbf{Backup Strategies:} Maintain passive gravity system as emergency fallback delivering >20\,gpm; integrate seamless switch‑over logic.
%   \item \textbf{Alternative Concepts:} Investigate modular pump stations scalable by flow requirement; evaluate micro‑hydro turbine integration for continuous recharge.
%   \item \textbf{Intellectual Property:} Pursue patenting of quick‑connect irrigation manifold and automated dual‑mode switching algorithm.
% \end{itemize}

% % 15.4 Lessons Learned
% \section*{15.4 Lessons Learned}
% \begin{itemize}
%   \item Phased development (scoping, infrastructure prep, final integration) accelerated risk mitigation and stakeholder engagement.
%   \item FMEA and House of Quality analysis directed critical trade‑offs between flow performance and energy consumption.
%   \item Python‑based data analysis facilitated rapid validation of system metrics against targets.
%   \item Challenges: component lead times, balancing aesthetic requirements with robust enclosure design, and incomplete long‑term durability data.
%   \item Future projects will benefit from early prototype reviews, cross‑disciplinary site visits, and clearer test plans for field fatigue testing.
% \end{itemize}

% \end{document}



\documentclass[12pt]{report}
\usepackage[table]{xcolor}
\usepackage[margin=1in]{geometry}
\usepackage{booktabs}
\usepackage{hyperref}
\usepackage{array}
\usepackage{float}
\usepackage{caption}
\usepackage{longtable}
\geometry{margin=1in}
\begin{document}

\chapter{Conclusions}

% 15.1 Summary of Design Project Results
\section{Summary of Design Project Results}
This design project, undertaken in collaboration with the USC Office of Sustainability and USC Facilities Department, focused on retrofitting and enhancing the existing 1,000‑gallon rainwater collection tank at the Green Quad permaculture garden—a 0.5‑acre, student-run organic garden behind a LEED‑certified building. The tank’s original flow of 1–2\,gpm was insufficient for operational needs.
\newline\newline
Performance was evaluated using Python-based statistical analysis:
\begin{itemize}
  \item \textbf{Flow rate (N=120):} Mean $\bar Q = 12.14\pm0.12\,$gpm meets fallback requirement (8–12 gpm).  
  \item \textbf{Outlet pressure (N=120):} Mean $\bar P = 32.24\pm0.57\,$psi satisfies fallback (20–40 psi).  
  \item \textbf{Relief valve (n=3):} Opening pressure $71.7\pm2.3\,$psi meets safety threshold (\textless 75 psi).  
  \item \textbf{Energy cost:} With pump draw $P=1314\,$W and daily runtime $t=270\,$s, annual energy $E=35.99\,$kWh gives operating cost \$4.32/yr (\textless \$50/yr target).  
\end{itemize}

\begin{table}[H]
  \centering
  \caption{Performance Metrics vs. Target Specifications}
  \label{tab:performance}
  \begin{tabular}{lcccr}
    \toprule
    Parameter & Target & Fallback & Measured & Met? \\
    \midrule
    Flow rate      & 12--16\,gpm & 8--12\,gpm    & 12.14 $\pm$ 0.12\,gpm & Yes \\
    System pressure & 40--60\,psi & 20--40\,psi  & 32.24 $\pm$ 0.57\,psi & Yes (fallback) \\
    Energy cost    & \textless \$50/yr    & \textless \$150/yr     & \$4.32/yr              & Yes \\
    Durability     & —           & —             & Not tested             & N/A \\
    Ease of assembly & —         & —             & Not tested             & N/A \\
    Capital cost   & \textless \$2{,}000 & \textless \$2{,}500     & \$1,550                & Yes \\
    Safety         & GFCI, \textless 75 psi & GFCI, \textless 80 psi & Pass           & Yes \\
    \bottomrule
  \end{tabular}
\end{table}

\begin{table}[H]
  \centering
  \caption{Customer Needs Satisfaction}
  \label{tab:needs}
  \begin{tabular}{p{0.5\linewidth}c}
    \toprule
    \textbf{Customer Need} & \textbf{Met?} \\
    \midrule
    Time-efficient water transfer   & Yes \\
    Minimal manual intervention     & Yes \\
    Efficient water distribution    & Yes \\
    Continuous operation            & Yes \\
    Energy cost savings             & Yes \\
    Safety compliance               & Yes \\
    Durability                      & N/A \\
    Ease of assembly                & Yes \\
    Capital cost limit              & Yes \\
    \bottomrule
  \end{tabular}
\end{table}

% 15.2 Standards Compliance
\begin{longtable}{%
    p{1.2cm}   % Agency
    p{4cm}   % Standard
    p{6cm}   % Product Feature
    p{3.8cm}   % Link
}
\caption{Applicable Standards and Compliance\label{tab:standards}}\\

% --- first head (on page 1) ---
\toprule
\textbf{Agency} & \textbf{Standard}
  & \textbf{Product Feature} & \textbf{Link} \\
\midrule
\endfirsthead

% --- head for subsequent pages ---
\multicolumn{4}{l}{\small\itshape Table \thetable\ continued from previous page}\\
\toprule
\textbf{Agency} & \textbf{Standard}
  & \textbf{Product Feature} & \textbf{Link} \\
\midrule
\endhead

% --- foot for pages before the last ---
\midrule
\multicolumn{4}{r}{\small\itshape Continued on next page\ldots}\\
\endfoot

% --- footer for the last page ---
\bottomrule
\endlastfoot

% --- now your data rows ---
OSHA  & Section 5(a)(1) of the OSH Act
      & GFCI‑protected circuit and lockable, weatherproof enclosure
      & \url{https://www.osha.gov/} \\

ASTM  & ASTM D1785‑19 – Standard Specification for PVC Plastic Pipe, Schedules 40, 80, and 120
      & Uses ¾" Schedule 40 PVC pipe manufactured to ASTM D1785 tolerances
      & \url{https://store.astm.org/} \\

ANSI  & NSF/ANSI 61 – Drinking Water System Components – Health Effects
      & NSF/ANSI 61 certified water-contact materials
      & \url{https://webstore.ansi.org/} \\

UL    & UL 778 – Standard for Motor-Operated Water Pumps
      & UL-listed motor-operated water pump
      & \url{https://www.shopulstandards.com/} \\

NEMA  & NEMA 3R (per NEMA 250) – Enclosures for Outdoor Electrical Equipment
      & NEMA 3R-rated enclosure
      & \url{https://www.nemaenclosures.com/} \\

ISO   & ISO 12100:2010 – Safety of Machinery
      & FMEA conducted per ISO 12100 methodology
      & \url{https://www.iso.org/} \\

ASME  & ASME B1.20.1:2013 – Pipe Threads, General Purpose (Inch)
      & NPT fittings machined to ASME B1.20.1 dimensions
      & \url{https://www.asme.org/} \\

ANSI  & ANSI Z535.4‑2023 – Product Safety Signs and Labels
      & Safety labels designed per ANSI Z535.4 specifications
      & \url{https://blog.ansi.org/} \\

NIST  & NIST Fluid Metrology’s Liquid Flow Standards
      & Flow‑metering devices calibrated to NIST gravimetric standards
      & \url{https://www.nist.gov/} \\

\end{longtable}

% 15.3 Recommendations for Future Development
\section{Recommendations for Future Development}
\begin{itemize}
  \item \textbf{Next Steps:} Conduct accelerated lifetime and field fatigue tests for tank integrity and valve wear; finalize supply chain and cost analysis for commercial scaling.
  \item \textbf{Commercialization Readiness:} Prototype meets performance and safety targets; refine manufacturing for enclosure and brass components; develop user documentation and installation guidelines.
  \item \textbf{Design Refinements:} Optimize UV‑resistant enclosure materials; implement snap‑fit modular connectors for tool‑free assembly; calibrate relief‑valve pre‑set for tighter tolerance.
  \item \textbf{Viability Assessment:} Core design is viable; if long‑term durability targets are not met, explore composite liners or secondary containment options.
  \item \textbf{Alternative Concepts:} Assess micro‑hydro turbine integration for energy recovery in larger installations; prototype battery‑less solar pumping modules for remote sites.
  \item \textbf{Intellectual Property:} No Intellectual Property worth pursuing was developed during this project's duration.
\end{itemize}

% 15.4 Lessons Learned
\section{Lessons Learned}
\begin{itemize}
  \item Phased prototyping and stakeholder workshops revealed operational insights and accelerated design validation.
  \item FMEA and House of Quality analyses guided prioritization, balancing flow performance, energy use, and safety.
  \item Python‑based statistical scripts provided rigorous validation of metrics but highlighted need for standardized test fixtures.
  \item Challenges included component lead times, UV exposure on enclosures, and limited long‑term durability data.
  \item Future projects should emphasize early supplier engagement, comprehensive field testing plans, and iterative user feedback cycles.
\end{itemize}

\end{document}
